\documentclass[tikz,border=10pt]{standalone}
% By Jan Suchanek 2022. CC-BY-2.5
% Phase oriented project management.
% Following: Schliersmann, Christiane and Heinz Ulrich Thiel:
% Organisationsentwicklung. Prinzipien und Strategien von
% Veränderungsprozessen. 5th edition, Springer, 2018, page 50.
\usepackage{tikz}
\begin{document}
     \def\radi{4cm}\relax  % overall radius
     \def\radii{1cm}\relax % radius of nodes
     \begin{tikzpicture}[rotate=90, every node/.style={font=\sffamily}]
         \draw (0:\radi) node {
             \parbox{1.5\radii}{\centering
                 analyze situation
             }
         } circle (\radii) ;
         \draw (-60:\radi) node {
             \parbox{1.5\radii}{\centering
                 clarify goals
             }
         } circle (\radii) ;
         \draw (-120:\radi) node {
             \parbox{1.5\radii}{\centering
                 work breakdown
             }
         } circle (\radii) ;
         \draw (-180:\radi) node {
             \parbox{1.5\radii}{\centering
                 planning
             }
         } circle (\radii) ;
         \draw (-240:\radi) node {
             \parbox{1.5\radii}{\centering
                 implement \& control
             }
         } circle (\radii) ;
         \draw (-300:\radi) node {
             \parbox{1.5\radii}{\centering
                 evaluate \& transfer
             }
         } circle (\radii) ;
         \foreach \i in {0,60,...,360}
             {
                 \draw[<-,thick] (\i+20:\radi) arc (\i+20:\i+40:\radi);
                 \foreach \ii in {0,60,...,360}
                     \draw  (\i:\radi-\radii) -- (\ii:\radi-\radii);
             }
     \end{tikzpicture}
\end{document}
